\chapter{Introduction: Logic and Computation}

\section{Why study logic in computer science?}

Logic is the study of \emph{correct reasoning}: which conclusions follow from which
assumptions, independent of what the statements are actually about. CS 245 treats logic
as a mathematical object in its own right. We define precise \emph{languages}
(propositional logic and first-order predicate logic), give them a precise
\emph{meaning} (semantics), and design \emph{proof systems} that let us derive conclusions
by symbol manipulation alone. We then ask how the two sides relate (soundness and
completeness) and what the fundamental limits of the approach are (undecidability,
incompleteness).

Logic shows up all over computer science:
\begin{itemize}
  \item \textbf{Hardware}: digital circuits are propositional formulas; circuit
        minimisation is simplification of Boolean formulas.
  \item \textbf{Programming languages}: type systems are logics (the Curry--Howard
        correspondence: proofs are programs, propositions are types); conditionals and
        loop guards are formulas.
  \item \textbf{Program verification}: specifications are formulas, and correctness
        proofs (Hoare logic, Chapter~\ref{ch:hoare}) are formal deductions.
  \item \textbf{Databases}: SQL queries are essentially first-order formulas (relational
        calculus).
  \item \textbf{Artificial intelligence and automated reasoning}: SAT solvers, SMT
        solvers, theorem provers, model checkers, Prolog.
  \item \textbf{Theory of computation}: the question ``is there an algorithm that decides
        whether a formula is valid?'' leads directly to Turing machines and the halting
        problem (Chapter~\ref{ch:undecidability}).
\end{itemize}

\section{Arguments, premises and conclusions}

\begin{definition}[Proposition]
A \emph{proposition} is a declarative sentence that is either true or false (but not
both).
\end{definition}

\begin{example}
``$2+2=4$'', ``Waterloo is in Ontario'' and ``every even integer greater than $2$ is the sum
of two primes'' are propositions (the last one has an unknown truth value, but it has
one). ``Close the door!'', ``Is it raining?'' and ``$x > 3$'' (with $x$ unspecified) are
\emph{not} propositions.
\end{example}

\begin{definition}[Argument]
An \emph{argument} consists of a set of propositions called \emph{premises} and one
proposition called the \emph{conclusion}. The argument is \emph{valid} if, in every
situation in which all the premises are true, the conclusion is also true.
\end{definition}

The key word is \emph{every}. Validity depends only on the \emph{form} of the argument,
not on whether the premises are actually true.

\begin{example}[A valid argument]
\begin{quote}
If it is raining, then the ground is wet.\\
It is raining.\\
Therefore, the ground is wet.
\end{quote}
Writing $p$ for ``it is raining'' and $q$ for ``the ground is wet'', the form is
\[
  p \to q,\; p \;\;\therefore\;\; q.
\]
This form (\emph{modus ponens}) is valid no matter what $p$ and $q$ stand for.
\end{example}

\begin{example}[An invalid argument]
\begin{quote}
If it is raining, then the ground is wet.\\
The ground is wet.\\
Therefore, it is raining.
\end{quote}
The form $p\to q,\ q \;\therefore\; p$ (\emph{affirming the consequent}) is invalid: take
$p$ false and $q$ true (someone ran the sprinkler). Both premises are true and the
conclusion is false. One such situation, a \emph{counterexample}, is enough to show
invalidity.
\end{example}

\begin{important}[Valid vs.\ sound vs.\ true]
\begin{itemize}
  \item A \emph{valid} argument can have false premises and a false conclusion:
        ``All cats are fish; all fish fly; therefore all cats fly'' is valid.
  \item An argument is \emph{sound} (in the informal sense) if it is valid \emph{and} all
        its premises are actually true. Then its conclusion must be true.
  \item An invalid argument can have a true conclusion; that the conclusion happens to be
        true does not make the reasoning correct.
\end{itemize}
(Careful: later we use ``soundness'' for a property of a \emph{proof system}, which is a
different idea.)
\end{important}

\section{Two views of ``follows from''}

CS 245 studies two quite different ways of making ``the conclusion follows from the
premises'' precise.

\begin{description}
  \item[Semantic (model-theoretic).] Consider every possible way of assigning meaning to
        the symbols (every truth valuation, or every interpretation in predicate logic).
        If the conclusion is true in every one in which all premises are true, we write
        $\Sigma \models \varphi$ (\emph{$\Sigma$ semantically entails $\varphi$}).
  \item[Syntactic (proof-theoretic).] Fix a finite set of inference rules. If there is a
        finite derivation of $\varphi$ from $\Sigma$ using only those rules, we write
        $\Sigma \vdash \varphi$ (\emph{$\varphi$ is provable from $\Sigma$}).
\end{description}

The semantic view is close to our intuition about ``truth'', but checking infinitely many
interpretations is not mechanical. The syntactic view is mechanical (a computer can check
a proof), but on its face has nothing to do with truth. The central results of the course
are that, for propositional and first-order logic, the two coincide:
\[
  \Sigma \vdash \varphi \iff \Sigma \models \varphi
  \qquad\text{(soundness $\Rightarrow$, completeness $\Leftarrow$).}
\]

\section{Formal languages: object language vs.\ metalanguage}

When we study a logic we use two languages at once:
\begin{itemize}
  \item the \emph{object language}: the formal logic being studied, with symbols such as
        $p, q, \neg, \land, \to, \forall$;
  \item the \emph{metalanguage}: the (mathematical) English we use to talk \emph{about}
        the object language, including symbols like $\models$, $\vdash$, ``iff'' and
        $\Rightarrow$.
\end{itemize}
For example, $p \to q$ is a formula of the object language, while
``$\{p\to q, p\} \models q$'' is a statement in the metalanguage that happens to be true.
Keeping the two levels apart avoids a lot of confusion, especially when comparing $\to$
(a connective inside formulas) with $\Rightarrow$ or ``implies'' (reasoning about
formulas).

\section{Translating English into logic}

A recurring skill in the course is translating English sentences into formulas. Some
standard translations:

\begin{center}
\begin{tabular}{p{0.62\textwidth}l}
\toprule
English & Formula \\
\midrule
not $p$; it is not the case that $p$ & $\neg p$ \\
$p$ and $q$; $p$ but $q$; $p$ although $q$ & $p \land q$ \\
$p$ or $q$ (inclusive) & $p \lor q$ \\
if $p$ then $q$; $q$ if $p$; $p$ only if $q$; $p$ is sufficient for $q$;
  $q$ is necessary for $p$ & $p \to q$ \\
$p$ if and only if $q$; $p$ exactly when $q$ & $p \leftrightarrow q$ \\
$p$ unless $q$ & $\neg q \to p$ (equivalently $p \lor q$) \\
either $p$ or $q$ but not both & $(p\lor q)\land\neg(p\land q)$ \\
neither $p$ nor $q$ & $\neg p \land \neg q$ \\
\bottomrule
\end{tabular}
\end{center}

\begin{important}[Pitfall: ``only if'']
``$p$ only if $q$'' means $p \to q$, \emph{not} $q \to p$. ``You can graduate only if you
pass CS 245'' says graduating requires passing; it does not say passing guarantees
graduation.
\end{important}

\begin{example}
``If Alice goes to the party and Bob does not, then Carol will be unhappy unless Dave
comes.'' With atoms $a,b,c,d$:
\[
  (a \land \neg b) \to (\neg d \to c).
\]
\end{example}

\section{Roadmap of the course}

\begin{enumerate}
  \item Propositional logic: syntax, semantics, proof systems, soundness and completeness.
  \item Predicate (first-order) logic: syntax, semantics, natural deduction, soundness
        and completeness.
  \item Applications: formal arithmetic (Peano axioms), program verification with Hoare
        logic.
  \item Limits: Turing machines, the halting problem, undecidability of first-order
        validity, G\"odel's incompleteness theorems.
\end{enumerate}

\section*{Summary}
\begin{itemize}
  \item A proposition is a declarative sentence with a definite truth value.
  \item An argument is valid iff every situation making all premises true makes the
        conclusion true; one counterexample shows invalidity.
  \item Validity is about form, not content; valid $\neq$ true.
  \item $\models$ (semantic entailment) and $\vdash$ (provability) are the two notions of
        consequence; soundness and completeness say they agree.
  \item Distinguish the object language from the metalanguage, and be careful
        translating ``only if'' and ``unless''.
\end{itemize}
